\chapter{Mathematical Demonstrations}
\label{app:mathematical_demonstrations}
This appendix shows that FELICS reduces to the metric of Matton et al.~\cite{mattonWalkTalkMeasuring2024} when every candidate interaction group contains a single concept. It outlines the design rationale of the metric rather than serving as a comprehensive proof.

\section*{Equality of Causal Effect between FELICS singleton-groups and Original Metric}
\begin{math_box}[label=ma:one_item_group]
    Let the candidate interaction group consist of exactly one concept,
    \begin{equation}
        s=\{C_m\}.
    \end{equation}

    The payout coefficient introduced in Section~\ref{sec:shapley_correction} becomes

    \begin{equation}
        \gamma(\mathbf{x},s)
        =
        \frac{v_{\mathbf{x}}(s)}
        {\sum_{C_i\in s}v_{\mathbf{x}}(\{C_i\})}
        =
        \frac{v_{\mathbf{x}}(\{C_m\})}
        {v_{\mathbf{x}}(\{C_m\})}
        =1.
    \end{equation}

    Since the reference set used to compute the correction factor
    $\bar\gamma(\mathbf{x},s,T)$ is restricted to sets of identical treatment size,
    the only admissible reference size is likewise $|T|=1$. Consequently,

    \begin{equation}
        \bar\gamma(\mathbf{x},s,T)=1,
    \end{equation}

    and the normalized value function reduces to

    \begin{equation}
        v_{\mathrm{adj},\mathbf{x}}(T)
        =
        \frac{v_{\mathbf{x}}(T)}
        {\bar\gamma(\mathbf{x},s,T)}
        =
        v_{\mathbf{x}}(T).
    \end{equation}

    The adjusted causal effect is then obtained from the Shapley formulation

    \begin{equation}
        \mathrm{CE}_{\mathrm{adj}}(\mathbf{x},C_m,s)
        =
        \sum_{T\subseteq s\setminus\{C_m\}}
        \frac{|T|!(|s|-|T|-1)!}{|s|!}
        \left[
            v_{\mathrm{adj},\mathbf{x}}(T\cup\{C_m\})
            -
            v_{\mathrm{adj},\mathbf{x}}(T)
            \right].
    \end{equation}

    Since $s=\{C_m\}$, the only admissible coalition is the empty set,
    $T=\emptyset$. The weighting coefficient therefore simplifies to

    \begin{equation}
        \frac{0!(1-0-1)!}{1!}=1,
    \end{equation}

    and using the convention $v_{\mathrm{adj},\mathbf{x}}(\emptyset)=0$,

    \begin{equation}
        \mathrm{CE}_{\mathrm{adj}}(\mathbf{x},C_m,s)
        =
        v_{\mathrm{adj},\mathbf{x}}(\{C_m\}).
    \end{equation}

    Because no normalization is applied for singleton groups,

    \begin{equation}
        v_{\mathrm{adj},\mathbf{x}}(\{C_m\})
        =
        v_{\mathbf{x}}(\{C_m\}),
    \end{equation}

    which is exactly the value function underlying the causal-effect definition of Matton et al.~\cite{mattonWalkTalkMeasuring2024}. Therefore,

    \begin{equation}
        \mathrm{CE}_{\mathrm{adj}}(\mathbf{x},C_m,\{C_m\})
        =
        \mathrm{CE}_{\mathrm{Matton}}(\mathbf{x},C_m)
    \end{equation}

    showing that FELICS reduces exactly to the original CE when every candidate interaction group contains only a single concept. Because the definitions of the explanation-implied effect and faithfulness (see Equations~\ref{eq:ee} and \ref{eq:faithfulness}) are unchanged, both metrics are equivalent in this scenario.
\end{math_box}
